\section{Introduction}
\label{intro:section}

For integers $s,t\ge2$, the Ramsey number $r(s,t)$ is the least integer
$N$ such that every finite simple graph on $N$ vertices contains a clique
of order $s$ or an independent set of order $t$. We study the
off-diagonal regime in which $s$ is fixed and $t$ tends to infinity
through the integers. All logarithms are natural.

\begin{theorem}\label{intro:main}
For every fixed integer $s\ge6$ there is a constant $C_s>0$ such that,
for every $\varepsilon>0$ and every sufficiently large integer $t$,
\[
 \frac{t^{s-1}}{(\log t)^{s-2+\varepsilon}}
 \le r(s,t)
 \le C_s\frac{t^{s-1}}{(\log t)^{s-2}}.
\]
In particular,
\[
 \lim_{t\to\infty}
 \frac{(s-1)\log t-\log r(s,t)}{\log\log t}=s-2.
\]
\end{theorem}

The theorem determines the logarithmic exponent. It does not assert
that $r(s,t)$ is bounded below by a positive constant times
$t^{s-1}/(\log t)^{s-2}$.

\subsection{Background and contribution}

Ramsey's finite partition theorem, developed in his work on a problem
of formal logic, guarantees that these numbers are finite
\cite[Theorem~B]{Ramsey1930}.
The classical argument of Erd\H{o}s and Szekeres gives
\[
 r(s,t)\le\binom{s+t-2}{s-1}
\]
and hence the polynomial upper bound $O_s(t^{s-1})$
\cite{ErdosSzekeres1935}. Ajtai, Koml\'os, and Szemer\'edi obtained the
logarithmic improvement
$O_s(t^{s-1}/(\log t)^{s-2})$ \cite{AKS1980}.
Li, Rousseau, and Zang subsequently proved the upper estimate with
leading coefficient $1+o(1)$ \cite{LiRousseauZang2001}.
The logarithmic power $s-2$ thus comes from the classical upper-bound
side of the problem. Determining whether a lower construction attains
that power requires controlling independent sets far more precisely
than the polynomial exponent alone.

For $s=3$, Kim's lower bound established
$r(3,t)=\Theta(t^2/\log t)$ \cite{Kim1995}.
For fixed $s\ge5$, Spencer's local-lemma construction gives
$r(s,t)=\Omega_s((t/\log t)^{(s+1)/2})$ \cite{Spencer1977}.
Bohman and Keevash improved the logarithmic factor through the random
$K_s$-free process, proving
\[
 r(s,t)=\Omega_s\!\left(
 t^{(s+1)/2}(\log t)^{1/(s-2)-(s+1)/2}\right)
\]
\cite[Theorem~1.2]{BohmanKeevash2010}.
A different route uses finite geometry and pseudorandom graphs.
Mubayi and Verstra\"ete showed that suitable density-optimal
$K_s$-free pseudorandom graphs would yield the lower order
$t^{s-1}/(\log t)^{2s-4}$ \cite[Corollary~2]{MubayiVerstraete2024}; this is a
conditional construction criterion.
For $s=4$, Mattheus and Verstra\"ete obtained the unconditional bound
$\Omega(t^3/(\log t)^4)$ using a finite-geometric construction
\cite{MattheusVerstraete2024}.
Brada\v{c}'s projective construction then established
\[
 r(s,t)\ge c_s\frac{t^{s-1}}{(\log t)^{2s-4}}
 \qquad(s\ge3,\ t\ge2)
\]
\cite[Theorem~1.1]{Bradac2026}, determining the polynomial exponent
for every fixed $s$. The remaining difference between the powers
$2s-4$ and $s-2$ is the logarithmic gap addressed here.

The geometric ingredients also have an earlier history.
Projective orthogonality and its spectral estimates appear in
Alon and Krivelevich \cite{AlonKrivelevich1997}.
Ordered incidence constructions of Codenotti, Pudl\'ak, and Resta
\cite{CodenottiPudlakResta2000}, and Kostochka, Pudl\'ak, and R\"odl
\cite{KostochkaPudlakRodl2010}, are predecessors of the flag approach.
We use the ordered incident-flag graph of Brada\v{c}'s
higher-dimensional construction \cite[Section~2.5]{Bradac2026}
and adapt its marking argument \cite[Claim~2.13]{Bradac2026}.
The distinction between a step that substantially shrinks a candidate
set and a step with few possible choices also appears in the
independent-set counting argument of Alon and R\"odl
\cite[proof of Theorem~2.1]{AlonRodl2005}. Here it must be combined with an entropy estimate
for a sequence selected from the random stream: selection need not
preserve the independence or uniformity of its flags.

The companion paper \emph{The sharp logarithmic exponent of $r(5,t)$}
\cite[Lemma~3.4 and Theorem~4.1]{RamseyFive} develops the selected-law entropy and
low-dimensional sparse-pair framework for $s=5$.
The present paper supplies all the required arguments locally.
To pass to arbitrary fixed dimension, we prove a repeated-projection
reduction for sparse pairs and a two-row high-rank description using
a rich-subspace bound that follows from Nie and Wang's finite-degree
closure inequality \cite{NieWang2015}.
The repeated projections and the high-rank construction require their
own estimates; neither is inferred from the five-clique theorem.
The low-dimensional polynomial argument adapts methods of Dvir,
Guth--Katz, and Elekes--Kaplan--Sharir
\cite{Dvir2009,GuthKatz2010,ElekesKaplanSharir2011}, with the
positive-characteristic obstructions analyzed by Ellenberg and
Hablicsek explicitly controlled \cite{EllenbergHablicsek2016}.

\subsection{The construction and its entropy obstruction}

The main work is the following prime-indexed construction. For a graph
$G$, write $\alpha(G)$ for its independence number.

\begin{theorem}\label{intro:prime-construction}
Fix an integer $d\ge5$ and a real number $0<\eta<1/10$.
For every sufficiently large prime $q$, with $\sigma=\log q$, there is
a graph $G$ such that
\[
 |V(G)|=\lfloor q^d\sigma\rfloor,\qquad
 K_{d+1}\not\subseteq G,\qquad
 \alpha(G)<\lfloor q\sigma^{1+\eta}\rfloor.
\]
\end{theorem}

Here is the construction and the organization of its proof. A point of
$\PG(d,q)$ is a one-dimensional vector subspace of $\F_q^{d+1}$.
A dual point is a one-dimensional subspace of the dual vector space;
it specifies a projective hyperplane. Write $a\perp b$ when the
covector $b$ vanishes on the point $a$. An incident flag is a pair
$(a,b)$ with $a\perp b$.

Choose $N=\lfloor q^d\sigma\rfloor$ flags independently and uniformly,
and use their positions as vertices. Positions $i<j$ are adjacent when
\[
 a_i\perp b_j\quad\hbox{and}\quad a_j\not\perp b_i.
\]
Linear independence excludes $K_{d+1}$. An independent set, in its
original order, is therefore a flag sequence satisfying
\[
 i<j,\quad a_i\perp b_j\quad\Longrightarrow\quad a_j\perp b_i.
\]
We call this condition consistency. Its asymmetry records the order;
reversing the order and interchanging the two endpoint spaces preserves
the condition.

Suppose that every stream of flags contained a consistent sequence of
length $k=\lfloor q\sigma^{1+\eta}\rfloor$. Selecting such a sequence
can heavily bias its distribution. Nevertheless, a counting argument
gives every selected sequence $F$ of deterministic length
$\ell=\Theta(k)$ the entropy lower bound
\[
 H(F)\ge d\sigma\ell+\eta\ell\log\sigma-O(k).
\]
This bound survives any fixed number of restrictions to events of
probability bounded below. We will contradict it by successively
describing shorter sequences, still of length $\Theta(k)$, in smaller
flag domains.

The cost of the exceptional positions in the marking argument explains
why these successive descriptions are needed. Brada\v{c}'s direct count
records $O(q\sigma)$ expensive flags, with $O(\sigma)$ information
per flag \cite[proof of Lemma~2.12]{Bradac2026}. The resulting
$O(q\sigma^2)$ allowance is too coarse for the present entropy
comparison: its scale exceeds the surplus $\eta\ell\log\sigma$ when
$\ell=\Theta(q\sigma^{1+\eta})$. The leading term $d\sigma\ell$
already allows $d\sigma$ units for each retained flag. Smaller flag
domains reduce the excess cost of the expensive positions beyond this
allowance. At each stage we replace the old description instead of
carrying its history into the new message.

The description has three main components.

\paragraph{Describing a sparse pair of point sets.}
Given point sets $S,T$ in opposite point systems of $\PG(j,q)$, where
$2\le j\le d$, assume that
$|S||T|$ is close on a logarithmic scale to $q^{j+1}$, while their
incidence density is much smaller than $1/q$. The geometric description
lemma produces a short message specifying a set that captures a fixed
fraction of the smaller support and is not much larger than that
support. Separate validation samples then produce small domains on
both sides that retain almost all of $S$ and $T$.
Sections~\ref{geom:section} and~\ref{learn:section} prove this statement.
Repeated projection reduces dimensions $j\ge4$ to dimensions two and
three. In the remaining dimensions, the sampling argument uses independent
Poisson batches. A moment argument controls the size of the decoded set;
its geometric input is a rich-line estimate with the field characteristic
kept larger than the relevant polynomial degrees.

\paragraph{Finding supports to which the description applies.}
Consistency lets us mark all but $O(q\log q)$ positions as belonging
to explicitly known flag domains of size $O(q^d)$. The entropy deficit
in these domains bounds the dependence between suitable representative
positions and the remaining sequence. Most representative pairs
therefore have almost no one-way incidence conflict. In the reciprocal
classes, whose two prescribed endpoint-domain bounds have product
$O(q^{d+1})$, a geometric argument turns this into sparse incidence
between endpoint supports. Possible exact
incidence coincidences are charged to small occupied flag rectangles.
Section~\ref{mark:section} supplies these estimates. The other classes
have large endpoint ranks; Section~\ref{comp:section} describes them
directly using two independent sample rows and a polynomial bound for
the union of rich subspaces.

\paragraph{Iterating without retaining old information.}
For reciprocal supports, arrange the representative pairs in time
order and apply the geometric description along a balanced binary
tree. Each successful node restricts one endpoint domain for each
child. Fresh validation samples bound the loss from each restriction
before any later success is imposed. An additive potential bounds the
total message length. Crucially, the new decoder starts with the full
projective spaces and reads only the new message: it does not need the
previous context, hidden supports, or their conditional laws.
Section~\ref{comp:section} proves this one-step improvement, and
Section~\ref{end:section} iterates it a bounded number of times. The
resulting entropy upper bound is $d\sigma\ell+O(k)$, contradicting the
displayed lower bound.

The tables used to select description samples are public auxiliary
randomness. We bound message length, not search time. After averaging
over these tables, they can be fixed before imposing the next
constant-probability restriction on the stream.

Section~\ref{base:section} establishes the projective and information
estimates used throughout. The final section also transfers
Theorem~\ref{intro:prime-construction} from prime parameters to all
sufficiently large integers $t$ and proves the upper bound in
Theorem~\ref{intro:main}.
